% !TeX root = ../main.tex
% Add the above to each chapter to make compiling the PDF easier in some editors.

\chapter{Introduction}\label{chapter:introduction}

Recovering three-dimensional scene structure from video is a central problem in computer vision, with applications ranging from novel view synthesis and augmented reality to autonomous navigation and robotics. A prerequisite for virtually all of these tasks is knowledge of the camera geometry: both the intrinsic parameters that describe the camera's internal optics and the extrinsic parameters that describe how the camera moves through space. While controlled laboratory settings can provide these parameters through calibration targets or specialised hardware, the vast majority of available video data---from consumer devices, dashcams, surveillance feeds, and online platforms---is captured with unknown cameras under uncontrolled conditions.

Estimating camera intrinsics and extrinsics from such casual, in-the-wild video presents a significant challenge. The scenes are frequently dynamic, containing moving objects, people, and non-rigid content that violate the static-world assumptions of classical multi-view geometry. Camera motion is arbitrary and unpredictable, and intrinsic parameters vary between devices and recording conditions. Ground-truth annotations for camera parameters are expensive or infeasible to obtain at scale, limiting the applicability of supervised learning approaches.

\section{Motivation and Problem Statement}\label{sec:motivation}

Classical Structure-from-Motion (SfM) systems such as COLMAP~\parencite{schonberger2016structure} estimate camera parameters through iterative feature matching and bundle adjustment. While these methods achieve strong multi-frame consistency through global optimisation, they require test-time iteration, assume a predominantly static scene, and can fail on casual videos with weak texture, dynamic content, or poor camera trajectories. Supervised learning-based approaches~\parencite{wang2024dust3r,yugay2025vot,lin2025depth} replace test-time optimisation with feed-forward inference, improving speed and robustness, but they require ground-truth pose or geometry labels for training. This supervision is expensive to acquire at scale and limits the diversity of training data, constraining generalisation to new domains and camera types.

Recent progress in self-supervised learning has produced strong pretrained models for individual sub-tasks: monocular depth estimation, optical flow prediction, and single-image camera calibration can each be performed reliably by specialised networks. However, these models operate in isolation---a depth estimator has no awareness of camera motion, a calibration network has no temporal consistency, and a pose predictor may handle intrinsics through crude approximations. The question that motivates this thesis is: \emph{can we fuse multiple pretrained models into a unified system that jointly reasons about calibration and pose, while preserving the self-supervised nature of training?}

Self-supervised pose estimation methods such as AnyCam~\parencite{wimbauer2025anycam} demonstrate that camera poses can be learned from unlabelled video using flow reprojection losses, without any ground-truth annotations. However, existing self-supervised systems typically handle calibration through coarse approximations---for instance, AnyCam selects among 32 predefined focal length candidates, a procedure that is both computationally expensive and fundamentally limited in accuracy. Meanwhile, dedicated calibration networks like AnyCalib~\parencite{tirado2025anycalib} can predict accurate intrinsics from a single image by regressing per-pixel ray directions, but they process each frame independently with no temporal consistency.

These observations suggest a natural decomposition: use a dedicated calibration network for accurate intrinsics, and fuse its predictions into a self-supervised pose estimator. The key challenge is that per-frame calibration predictions are noisy---a single camera has fixed intrinsics throughout a video, yet independent predictions vary from frame to frame. Enforcing multi-frame consistency at the feature level, before the calibration network decodes its predictions, should produce more accurate and stable results.

\section{Challenges}\label{sec:challenges}

Three challenges arise when fusing pretrained models for joint calibration and pose estimation:

\paragraph{Multi-frame calibration consistency.}
A dedicated calibration network processes each frame independently, producing noisy and inconsistent predictions across a video sequence. Since camera intrinsics are constant within a sequence, these per-frame estimates should be aggregated. Simple scalar averaging discards the rich spatial structure in the network's intermediate features. An aggregation mechanism that operates at the feature level is needed.

\paragraph{Calibration--pose coupling.}
Calibration and pose are coupled through the projection equation: incorrect intrinsics lead to incorrect pose estimates, and vice versa. For small camera displacements, the intrinsic matrix cancels out of the reprojection equation entirely, leaving calibration unconstrained by the self-supervised loss. A stabilisation mechanism is needed to prevent calibration drift during training.

\paragraph{Efficient multi-frame training.}
Enforcing consistency across $N$ frames naively requires $O(N^2)$ pairwise flow computations. An efficient mechanism for multi-frame training that scales linearly is needed.

\section{Contributions}\label{sec:contributions}

This thesis proposes a framework for fusing a pretrained calibration network with a self-supervised pose estimator, introducing multi-frame calibration consistency via a learned feature aggregation module. Two contributions are made:

\begin{enumerate}
    \item \textbf{Multi-Frame Calibration Transformer (MCT).} A transformer module that aggregates spatial features from a pretrained calibration backbone across $N$ video frames before decoding, enforcing calibration consistency at the feature level. The MCT operates via cross-frame self-attention at each spatial position, then mean-pools across frames to produce a single consistent feature representation. It adds approximately 25M trainable parameters while all pretrained components remain frozen (${\sim}7.5\%$ of total parameters trained). The MCT is architecture-agnostic: it can be inserted between any feature extractor and decoder that share compatible spatial feature maps.

    \item \textbf{Demonstration on AnyCam and AnyCalib.} The framework is instantiated by fusing AnyCam (a self-supervised pose estimator) with AnyCalib (a single-image calibration network). The MCT is inserted between AnyCalib's frozen DINOv2 ViT-L/14 backbone and its Light-DPT decoder, and the resulting calibration is injected into AnyCam's pose head via a learned focal embedding. The fused system improves both pose accuracy (up to 39.5\% translation error reduction) and calibration accuracy (up to 32.5\% MAPE reduction) over the individual baselines, demonstrating that the framework effectively leverages complementary pretrained models.
\end{enumerate}

Training is fully self-supervised, combining flow reprojection losses with a calibration anchor that stabilises intrinsics during training. Multi-frame consistency is enforced efficiently through flow composition, reducing the cost from $O(N^2)$ to $O(N)$. A staged training pipeline ensures stable convergence.

\section{Thesis Outline}\label{sec:outline}

The remainder of this thesis is organised as follows.

Chapter~\ref{chapter:notation_guide} establishes the mathematical notation and symbol conventions used throughout the thesis. Chapter~\ref{chapter:fundamentals} introduces the geometric foundations---camera projection, depth, optical flow, multi-view consistency, and self-supervised learning---as well as the neural network architectures (Vision Transformers, DPT) that underpin modern geometric vision systems.

Chapter~\ref{chapter:related-work} surveys related work in camera pose and intrinsic estimation, comparing optimisation-based, supervised, and self-supervised approaches as well as standalone helper methods for depth, calibration, and optical flow.

Chapter~\ref{chapter:solution-approach} presents the technical solution: the Multi-Frame Calibration Transformer, its integration into a fused calibration--pose pipeline, and the self-supervised training procedure. The framework is described in generic terms and then instantiated with AnyCam and AnyCalib.

Chapter~\ref{chapter:evaluation} describes the evaluation methodology and experimental results, demonstrating the effectiveness of the proposed approach.

Chapter~\ref{chapter:future-work} discusses directions for future research, and Chapter~\ref{chapter:conclusion} concludes the thesis.
